\section{Paradigma Greedy}
\subsection{Proprietà della scelta greedy}
\subsubsection*{Scelta greedy}
Il paradigma greedy è un caso particolare di divide and conquer, in cui la soluzione di un'istanza viene costruita ad ogni
passo induttivo, attraverso la scelta di un elemento ottimo locale, con l'obiettivo di ottenere una soluzione globale ottima.
La scelta effettuata è detta scelta greedy e spesso è estremamente banale (con costo costante), la parte complessa è
dimostrare che la scelta greedy porta ad una soluzione globale ottima.

\subsubsection*{Proprietà di scelta greedy}
La proprietà di scelta greedy è la condizione che deve soddisfare una scelta greedy per garantire che la scelta effettuata
localmente non comprometta l'ottimalità della soluzione globale. Una scelta greedy riduce sia l'istanza da risolvere che
lo spazio generale delle soluzioni (ottimali e subottimali). Per portare alla soluzione (e soddisfare la proprietà di scelta
greedy) è necessario che lo spazio delle soluzioni, dopo la riduzione, contenga ancora almeno una soluzione ottima globale. 

\subsubsection*{Dimostrazione della correttezza di un algoritmo greedy}
A livello pratico, per dimostrare che una scelta greedy soddisfa la proprietà di scelta greedy, è necessario dimostrare
che esiste una soluzione ottima globale che soddisfi la scelta greedy. Spesso si considera una generica soluzione
ottima globale e, se tale soluzione non contiene la scelta greedy, di dimostra che è possibile modificarla senza
comprometterne l'ottimalità, in modo che soddisfi la scelta greedy.

\subsubsection*{Algoritmi greedy subottimali}
Esistono alcuni problemi per cui sono stati sviluppati algoritmi greedy subottimali, che non garantiscono in assoluto il
raggiungimento della soluzione ottima globale, ma la raggiungono o ci si avvicinano molto con alta probabilità. Questi
algoritmi risultano essere molto più veloci degli algoritmi ottimi e sono impiegati in cui si è disposti a sacrificare
un po' di ottimalità in cambio di una maggiore velocità di esecuzione.

\subsection{Scheduling di attività}
\subsubsection*{Formalizzazione del problema}
Data una sequenza \(S=\left\{ a_1, a_2, \dots a_n \right\}\) di \(n\) attività \(a_i : [s_i, f_i)\), identificate dal proprio
istante di inizio \(s_i\) e di fine \(f_i\) con \(s_i < f_i\), si vuole individuare un sottoinsieme fattibile
\(A \subseteq S\) di dimensione massima.

\begin{itemize}[topsep=5pt]
	\item un sottoinsieme \(A \subseteq S\) di attività è fattibile se composto da attività mutualmente compatibili
	\item due attività \(a_i\) e \(a_j\) sono compatibili se e solo se i loro intervalli di tempo non si sovrappongono:
	\vspace{-5pt}
	\[a_i : [s_i, f_i) \cap a_j : [s_j, f_j) = \emptyset \qquad \Longleftrightarrow \qquad s_i \geq f_j \;\; \vee \;\; s_j \geq f_i\]
\end{itemize}

\subsubsection*{Definizione della scelta greedy}
Data una sequenza \(S\) di \(n\) attività, si costruisce l'insieme soluzione \(A\) scegliendo l'attività \(a_i \in S\) che
termina per prima (con \(f_i = \min_j f_j\)) unita all'insieme soluzione \(A_\text{sub}\) della sottoistanza \(S_\text{sub}\)
ottenuta partendo da \(S\) e rimuovendo tutte le attività incompatibili con l'attività scelta \(a_i\).

Di seguito si formalizza la scelta greedy attraverso la proprietà di sottostruttura espressa in termini di equazioni alle
ricorrenze.
\[A(S) = \begin{cases}
	\left\{ a_1 \right\} & \text{se } |S| = n = 1 \\[5pt]
	\displaystyle \left\{ \left\{ a_k \in S \text{ con } k = \argmin f_i \right\} \;\; \cup \;\; A( \, S - \left\{ a_l \in S \text{ incompatibili con } a_k \right\} \,) \right\} & \text{se } |S| = n > 1
\end{cases}\]

\subsubsection*{Algoritmo ricorsivo ottimizzato basato sulla scelta greedy}
Per rendere più pratica l'operazione di scelta dell'attività ottima, si ordina la sequenza \(S\) in ordine crescente rispetto
agli istanti di fine \(f_i\) delle varie attività, con costo di ordinamento \(\Theta(n \log_2 n)\). In questo modo l'attività
da scegliere sarà sempre la prima della sequenza \(S\) ordinata.

Inoltre si osserva che non è necessario visitare tutta la sequenza \(S\) per rimuovere le tutte le attività incompatibili
ogni volta che viene scelta un'attività \(a_i\). È sufficiente fermarsi non appena si trova la prima attività compatibile
\(a_j\) che sarà l'attività successiva da aggiungere all'insieme soluzione \(A\). Le attività incompatibili con \(a_i\)
saranno anche incompatibili con \(a_j\), per cui verranno rimosse successivamente senza influenzare la scelta ottima.

\begin{algo}{REC\_GREEDY\_SEL}{$\underline{s}$, $\underline{f}$, i}{$\underline{s}$ e $\underline{f}$ istanti di inizio e fine ordinati per $f_i$, indice $i$ della prima attività compatibile}{$A = \left\{ a_{i_1}, a_{i_2}, \dots a_{i_k} \right\} \subseteq S$ fattibile e di dimensione massima}
	\State \(n \gets \text{len}(\underline{s})\)
	\If{\(i > n\)} \Return \(\varnothing\) \Comment{caso base: non ci sono più attività da considerare} \EndIf
	\State \(A \gets \left\{ i \right\}\) \Comment{scelta greedy: si sceglie la prima attività della sequenza}
	\State \(j \gets i + 1\)
	\While{\(j \leq n\) \textbf{and} \(\underline{s}[j] < \underline{f}[i]\)} \(j \gets j + 1\) \Comment{si saltano le attività incompatibili con \(a_i\)} \EndWhile
	\State \Return \(A \cup \text{REC\_GREEDY\_SEL}(\underline{s}, \underline{f}, j)\) \Comment{proprietà di sottostruttura con chiamata ricorsiva}
\end{algo}

\subsubsection*{Dimostrazione della correttezza - verifica della proprietà greedy}
Di seguito si dimostra che la scelta greedy ideata rispetta la proprietà greedy, ovvero che esiste una soluzione ottima
\(A^*\) che rispetta la scelta greedy dell'algoritmo.

Si suppone di avere una generica soluzione ottima \(A*\):
\begin{itemize}
	\item se la prima attività \(a_1\) della soluzione ottima \(A^*\) è \(a_k\) con \(k = \argmin_i f_i\), allora la scelta
	greedy è banalmente rispettata
	\item se la prima attività \(a_1\) della soluzione ottima \(A^*\) non è \(a_k\) con \(k = \argmin_i f_i\), vuol dire che
	\(f_1 > f_k\), per cui è possibile sostituire \(a_1\) con \(a_k\) senza perdere fattibilità, ottenendo una nuova soluzione
	che rispetta la scelta greedy e rimane soluzione ottima
\end{itemize}

\subsubsection*{Analisi della complessità}
Escludendo il costo dell'ordinamento iniziale \(\Theta(n \log_2 n)\), la complessità dell'algoritmo è data soltanto dalle
operazioni di confronto durante la visita della sequenza \(S\) per saltare le attività incompatibili. Siccome le attività
della sequenza vengono visitate in maniera sequenziale e ad ogni visita vengono effettuate un numero costante di operazioni,
la complessità asintotica dell'algoritmo è \(T(n) = \Theta(n)\).

\newpage

\subsection{Codice di Huffman}
\subsubsection*{Introduzione al codice di Huffman}
Il codice di Huffman è un codice a prefisso utilizzato in teoria dell'informazione per codificare simboli con una certa
distribuzione di probabilità nota, in sequenze di bit dette parole di codice di lunghezza variabile. Ha la caratteristica
di essere ottimo, ovvero di minimizzare la lunghezza media dei codici generati.

Essendo un codice a prefisso, il codice di Huffman garantisce che nessuna parola di codice associata ad un simbolo sia
prefisso di un'altra parola associata ad un altro simbolo. Questo aspetto aiuta in fase di decodifica perché, appena si
individua una parola di codice associata ad un simbolo, la si può decodificare immediatamente e senza ambiguità.

Le prestazioni di un codice di Huffman sono misurate in termini di lunghezza media dei codici generati ed è definita come
la lunghezza media delle parole di codice, pesata in base alla loro probabilità:
\[L = \sum_{i=0}^{\# \; \text{parole}} p_i \cdot l_i \qquad \text{con} \qquad \begin{aligned}
	p_i & = \text{probabilità di occorrenza del simbolo } i \\[3pt]
	l_i & = \text{lunghezza della parola di codice associata al simbolo } i
\end{aligned}\]

\subsubsection*{Rappresentazione tramite Trie}
Il trie è un albero binario pieno (ogni nodo interno ha due figli), non completo (le foglie possono essere a livelli diversi)
utilizzato per rappresentare le parole di codice generate dal codice di Huffman. Ha le seguenti caratteristiche:
\begin{itemize}
	\item ogni nodo interno del trie ha esattamente due figli, quello di sinistra raggiungibile tramite un arco etichettato
	con un bit 0 e quello di destra raggiungibile tramite un arco etichettato con un bit 1
	\item ogni foglia del trie rappresenta un simbolo, la cui parola di codice è data dalla concatenazione dei bit etichettati
	sugli archi del percorso che va dalla radice alla foglia
\end{itemize}
Si osservano le seguenti proprietà del trie:
\begin{itemize}
	\item se una parola di codice fosse prefisso di un'altra, allora la parola più corta sarebbe rappresentata da un
	nodo interno del trie lungo il percorso che porta dalla radice alla foglia che rappresenta la parola più lunga; ciò è
	correttamente impossibile per come definito il trie
	\item si possono scambiare le due etichette \(0,1\) degli archi che collegano i due figli al padre senza alterare
	le prestazioni del codice
	\item si possono scambiare simboli che si trovano allo stesso livello del trie senza alterare le prestazioni del codice,
	anche se hanno nodi padre diversi
\end{itemize}

\subsubsection*{Algoritmo greedy e costruzione dell'albero}
La costruzione dell'albero avviene in maniera bottom-up:
\begin{itemize}
	\item si parte inizializzando una foresta di \(n\) alberi, ognuno costituito da un singolo nodo foglia che rappresenta
	un simbolo con la relativa probabilità di occorrenza
	\item ad ogni passo si scelgono i due alberi con le probabilità della radice più piccole, e si crea un nuovo nodo padre
	comune alle due radici, con probabilità pari alla somma delle probabilità delle due radici; per convenzione si collega
	l'albero con probabilità più piccola alla sinistra assegnando un bit 0 all'arco, e l'albero con probabilità più grande
	a destra assegnando un bit 1 all'arco
	\item si ripete il passo precedente fino a quando tutte le foreste non si sono unite in un unico albero, che rappresenta
	il trie del codice di Huffman e la radice ha probabilità pari a 1
\end{itemize}

La scelta greedy è quella di prendere i due alberi le cui radici hanno le probabilità più piccole.

\newpage

\subsubsection*{Implementazione greedy iterativa con priority queue}
\vspace{-7pt}
\begin{algo}{IT\_MCM}{$\underline{c}$, $\underline{f}$}{vettori $\underline{c}$ e $\underline{f}$ rispettivamente degli $n$ simboli e delle relative frequenze}{riferimento alla radice del trie che rappresenta il codice di Huffman}
	\State \(n \gets \text{len}(\underline{c})\)
	\State \(Q \gets \text{PriorityQueue}()\) \Comment{creazione di una coda di priorità vuota}
	\ForAll{\(c\)}{\(\underline{c}\)} \Comment{inizializzazione della coda di priorità con i nodi foglia}
		\State \(z \gets \text{new\_node}(\textit{symbol}: c, \;\; \textit{freq}: \underline{f}[c], \;\; \textit{left\_node}: \text{null}, \;\; \textit{right\_node}: \text{null})\)
		\State \(Q.\text{insert}(z.\text{freq}, z)\)
	\EndFor
	\For{\(i \gets 1\)}{\(n-1\)} \Comment{ad ogni iterazione diminuisco di 1 il numero di alberi nella foresta}
		\State \(x \gets Q.\text{extract\_min}() \;\; - \;\; y \gets Q.\text{extract\_min}()\)
		\State \(z \gets \text{new\_node}(\textit{symbol}: \text{null}, \;\; \textit{freq}: x.\text{freq} + y.\text{freq}, \;\; \textit{left\_node}: x, \;\; \textit{right\_node}: y)\)
		\State \(Q.\text{insert}(z.\text{freq}, z)\)
	\EndFor
	\State \Return \(Q.\text{extract\_min}()\)
\end{algo}
\vspace{-18pt}

\subsubsection*{Analisi della complessità computazionale}
Si sceglie di usare un heap binario per implementare la coda di priorità, che garantisce un costo \(\Theta(\log_2 n)\) per ogni
inserimento ed estrazione. La complessità dell'algoritmo è data da:
\begin{itemize}
	\item \(n + n-1 = 2n-1\) inserimenti nella coda di priorità, con costo \(\Theta(\log_2 n)\) per ogni inserimento
	\item \(2(n-1) + 1 = 2n-1\) estrazioni dalla coda di priorità, con costo \(\Theta(\log_2 n)\) per ogni estrazione
	\item costo lineare \(\Theta(n)\) dei due for, al netto delle operazioni di inserimento ed estrazione
\end{itemize}
\vspace{-7pt}

\[T(n) = (2n-1) \cdot T_{\, \text{insert()}} + (2n-1) \cdot T_{\, \text{extract\_min}()} + \Theta(n) = \Theta(n \log_2 n)\]

\noindent
Si osserva che, i primi \(n\) inserimenti delle foglie possono essere fatti in tempo \(\Theta(n)\) tramite l'algoritmo di
costruzione di un heap dato un array di \(n\) elementi. Siccome, però, l'algoritmo di Huffman richiede almeno \(n\) estrazioni
e non si dispone di algoritmi migliori di \(\Theta(\log_2 n)\) per l'estrazione, la complessità dell'algoritmo rimane \(\Theta(n \log_2 n)\).

\subsubsection*{Dimostrazione della correttezza}
Di seguito si dimostra che la scelta greedy di estrarre i due nodi di frequenza minima \(a\) e \(b\) per primi per
posizionali al livello più basso dell'albero, porta ad una soluzione ottima globale.
\begin{itemize}
	\item si considera un albero ottimo \(T\) con costo ottimo \(B(T) = \sum_{i \in T} f(i) \cdot d(i) \leq B(T^*)\) minore
	di qualsiasi altro albero \(T^*\) in cui:
	\begin{itemize}[topsep=-2pt]
		\item due foglie \(a\) e \(b\) che si trovano alla massima profondità dell'albero, uniti da un nodo padre \(p\)
		\item due foglie \(x\) e \(y\) con frequenze minime \(\max\left\{ f(x), f(y) \right\} \leq \min \left\{ f(a), f(b) \right\}\)
	\end{itemize}
	\item se \(a \equiv x\) e \(b \equiv y\), la scelta greedy è banalmente rispettata
	\item altrimenti si ipotizza di scambiare \(a\) con \(x\) e \(b\) con \(y\), ottenendo un nuovo albero \(T'\) e provocando
	una variazione del costo dell'albero \(\Delta B = B(T') - B(T)\) pari a:
	\begin{align*}
		\Delta B &= (f(a) \, d(x) \!+\! f(b) \, d(y) \!+\! f(x) \, d(a) \!+\! f(y) \, d(b)) - (f(a) \, d(a) \!+\! f(b) \, d(b) \!+\! f(x) \, d(x) \!+\! f(y) \, d(y)) \\
		&= f(a) \, (d(x) \!-\! d(a)) + f(y) \, (d(b) \!-\! d(y)) + f(x) \, (d(a) \!-\! d(x)) + f(b) \, (d(y) \!-\! d(b)) \\
		&= \underbrace{(f(a) - f(x))}_{\geq 0} \, \underbrace{(d(x) - d(a))}_{\leq 0} + \underbrace{(f(b) - f(y))}_{\geq 0} \, \underbrace{(d(y) - d(b))}_{\leq 0} \leq 0 \quad \Rightarrow \quad B(T') \leq B(T)
	\end{align*}
	\item da \(B(T) \leq B(T')\) per l'ipotesi iniziale e \(B(T') \leq B(T)\) per il punto sopra, si ha che \(B(T) = B(T')\),
	ovvero che anche l'albero \(T'\) che soddisfa la proprietà greedy è ottimo.
\end{itemize}

È utile notare che la scelta greedy non distingue se i nodi contenuti nella priority queue siano foglie o radici di
sottoalberi, ma si basa solo sulla loro frequenza. Ciò significa che la dimostrazione sopra (che lavora solo con le foglie)
vale anche per i nodi interni ad un certo livello dell'albero (che eventualmente potrebbero essere sostituiti da foglie
di pari frequenza senza modificare il costo dell'albero).

\newpage

\subsection{(2,1)-boxing}
\subsubsection*{Formalizzazione del problema}
Sia data una sequenza di oggetti \(O = \left\{ o_1, o_2, \dots o_n \right\}\) con la sequenza dei rispettivi pesi
\(P = \left\{ p_1, p_2, \dots p_n \right\}\) positivi e minori di 1,ordinati in peso crescente (\(0 < p_1 \leq p_2 \leq \dots \leq p_n \leq 1\)),
e sia data inoltre una sequenza di scatole \(B = \left\{ b_1, b_2, \dots b_m \right\}\), tali che ogni scatola
\(b_p = \{ o_i, o_j \}\) può contenere al massimo due oggetti \(o_i, o_j\) di peso massimo complessivo inferiore ad 1 (\(p_i + p_j \leq 1\)).

Il problema del (2,1)-boxing consiste nel trovare una disposizione degli oggetti \(O\) nelle scatole \(B\) che minimizzi
il numero di scatole utilizzate.

\subsubsection*{Definizione della scelta greedy}
Si sceglie di inserire per ogni scatola l'oggetto più pesante \(o_j\) non ancora assegnato (l'ultimo della sottostringa
\(O_{i:j}\) degli oggetti non ancora assegnati) con eventuale aggiunta del più leggero \(o_i\) non ancora assegnato (il
primo della sottostringa \(O_{i:j}\)) se il peso lo permette.

Di seguito si formalizza la scelta greedy attraverso la proprietà di sottostruttura espressa in termini di equazioni alle
ricorrenze.
\[B(O_{i:j}, \; P_{i:j}) = \begin{cases}
	\{b_p = \{o_i\}\} & \text{se } i = j \\[3pt]
	\{b_p = \{o_j\}\} \;\; \cup \;\; B(O_{i:j-1}, \; P_{i:j-1}) & \text{se } i < j \text{ e } p_i + p_j > 1 \\[3pt]
	\{b_p = \{o_i, \, o_j\}\} \;\; \cup \;\; B(O_{i+1:j-1}, \; P_{i+1:j-1}) & \text{se } i < j \text{ e } p_i + p_j \leq 1
\end{cases}\]

\subsubsection*{Dimostrazione della correttezza}
Per prima cosa si osserva che la scelta è fattibile, ovvero soddisfa il limite massimo di peso per ogni scatola imposto
dal problema in tutti e tre i casi della definizione ricorsiva della scelta greedy.

Si dimostra ora che esiste una soluzione ottima che soddisfa la scelta greedy.
\begin{itemize}
	\item si ipotizza una soluzione ottima \(B^*(O_{i:j}, \; P_{i:j})\) per l'istanza \(O_{i:j}\), \(P_{i:j}\) e si analizza
	la scatola contenente l'oggetto più pesante \(b_p = \{o_j, \dots\}\):
	\item se la scatola contiene solo l'oggetto più pesante \(b_p = \left\{o_j\right\}\) e vale \(p_i + p_j > 1\), allora
	la scelta greedy è banalmente rispettata
	\item se la scatola contiene solo l'oggetto più pesante \(b_p = \left\{o_j\right\}\), ma vale \(p_i + p_j \leq 1\),
	allora è possibile spostare l'oggetto più leggero \(o_i\) dalla scatola in cui è inserito alla scatola \(b_p\) senza
	violare il vincolo di peso massimo e senza modificare il numero di scatole utilizzate, ottenendo una nuova soluzione
	ottima che rispetta la scelta greedy
	\item se la scatola contiene sia l'oggetto più pesante che quello più leggero \(b_p = \left\{o_i, o_j\right\}\), allora
	la scelta greedy è banalmente rispettata siccome deve necessariamente valere \(p_i + p_j \leq 1\)
	\item se la scatola contiene l'oggetto più pesante e un oggetto che non è il più leggero \(b_p = \left\{ o_j, o_k \right\}\),
	e l'oggetto più leggero si trova in un'altra scatola \(b_q = \{o_i, o_h\}\) allora è possibile scambiare \(o_i\) con \(o_k\)
	senza violare il vincolo di peso massimo (come illustrato di seguito) e senza modificare il numero di scatole utilizzate,
	ottenendo una nuova soluzione ottima that rispetta la scelta greedy
	\[b_p = \{ o_j, o_k\} \;\Rightarrow\; \begin{cases}
		\; p_j + p_k \leq 1 \;\Rightarrow\; p_j + p_i \leq p_j + p_k \leq 1 & \text{per } p_i \leq p_k \text{ con } o_i \text{ il più leggero} \\
		\; p_j + p_k \leq 1 \;\Rightarrow\; p_h + p_k \leq p_j + p_k \leq 1 & \text{per } p_h \leq p_j \text{ con } o_j \text{ il più pesante}
	\end{cases}\]
\end{itemize}
